% Autorska priprava k Realisticky.cz 020210–020213; TP 4.11; ulohy jsou autorske a obmenene.
% Zamerne nenahrazuje vyklad presnymi kopiemi cizich textu/prikladu.
\documentclass[11pt,a4paper]{article}
\input{../../sablony/pripravy/priprava.sty}
\usepackage{tikz}
\nastavPripravu{1. ročník}{Rovnice a nerovnice}{93}{Vyjádření neznámé z technického vzorce}
\begin{document}
\blok{Návaznost a motivace}
Na SPŠ se často potřebujeme dostat k jiné veličině než té, která je původně vyjádřena. Využijeme ekvivalentní úpravy známé z lineárních rovnic, ale ohlídáme nenulové dělitele a fyzikální jednotky.
\blok{Princip a odvození}
Považuj ostatní písmena za známé parametry. Při dělení symbolem napiš podmínku, že není nulový. Po vyjádření veličiny udělej rozměrovou kontrolu, pokud jednotky dávají smysl.
\blok{Řešené příklady: od jednoduchého ke složitějšímu}
\begin{prikladReseniBox}[title={Příklad 1}]
Ze vztahu $s=vt$ vyjádřete $t$ a $v$.
\tcblower
\begin{align*}
s&=vt \\
 t&=\frac{s}{v}\quad(v\ne0) \\
 v&=\frac{s}{t}\quad(t\ne0)
\end{align*}

Podmínky patří k jednotlivým úpravám; nulová rychlost mění charakter modelu.
\end{prikladReseniBox}
\begin{prikladReseniBox}[title={Příklad 2}]
Z Ohmova zákona $U=RI$ vyjádřete $I$ a $R$.
\tcblower
\begin{align*}
U&=RI \\
 I&=\frac UR\quad(R\ne0) \\
 R&=\frac UI\quad(I\ne0)
\end{align*}
\end{prikladReseniBox}
\blok{Samostatné procvičování}
\begin{zadaniReseniBox}[title={Úloha 1}]
Z $P=2(a+b)$ vyjádřete $a$.
\tcblower
\begin{align*}
P&=2a+2b \\
 2a&=P-2b \\
 a&=\frac P2-b
\end{align*}
\end{zadaniReseniBox}
\begin{zadaniReseniBox}[title={Úloha 2}]
Z $A=\frac{(a+b)h}{2}$ vyjádřete $b$.
\tcblower
\begin{align*}
2A&=(a+b)h \\
 \frac{2A}{h}&=a+b\quad(h\ne0) \\
 b&=\frac{2A}{h}-a
\end{align*}
\end{zadaniReseniBox}
\begin{zadaniReseniBox}[title={Úloha 3}]
Z $F=ma$ vyjádřete $m$.
\tcblower
\begin{align*}
F&=ma \\
 m&=\frac Fa\quad(a\ne0)
\end{align*}
\end{zadaniReseniBox}
\blok{Volitelná rezerva}
\begin{zadaniReseniBox}[title={Úloha 1}]
Z $C=\frac{m_1c_1+m_2c_2}{m_1+m_2}$ vyjádřete $c_2$.
\tcblower
\begin{align*}
C(m_1+m_2)&=m_1c_1+m_2c_2 \\
 m_2c_2&=C(m_1+m_2)-m_1c_1 \\
 c_2&=\frac{C(m_1+m_2)-m_1c_1}{m_2}
\end{align*}

Podmínky: $m_1+m_2\ne0$ a $m_2\ne0$.
\end{zadaniReseniBox}
\blok{Domácí úkol}
\begin{zadaniReseniBox}[title={Úloha 1}]
Ze vztahu $E=Pt$ vyjádřete $t$.
\tcblower
\begin{align*}
E&=Pt \\
 t&=\frac EP\quad(P\ne0)
\end{align*}
\end{zadaniReseniBox}
\begin{zadaniReseniBox}[title={Úloha 2}]
Z $V=abc$ vyjádřete $c$.
\tcblower
\begin{align*}
V&=abc \\
 c&=\frac V{ab}\quad(ab\ne0)
\end{align*}
\end{zadaniReseniBox}
\begin{zadaniReseniBox}[title={Úloha 3}]
Ze vzorce $c=\frac{m}{V}$ vyjádřete $V$.
\tcblower
\begin{align*}
cV&=m \\
 V&=\frac mc\quad(c\ne0)
\end{align*}

Původní předpis předpokládá také $V\ne0$.
\end{zadaniReseniBox}
\poznamka{Orientační časové rozvržení: 5 + 12 + 12 + 10 + 6 = 45 min. Přednost má pochopení jádra; rezerva není součástí časového součtu. Domácí úlohy jsou určeny na samostatnou práci, řešení zůstávají učiteli.}
\end{document}
